\documentclass[11pt]{article}

\usepackage[margin=1in]{geometry}
\usepackage{microtype}
\usepackage{amsmath,amssymb,amsthm,mathtools}
\usepackage{enumitem}
\usepackage[hidelinks]{hyperref}
\usepackage[nameinlink,noabbrev]{cleveref}
\usepackage{xurl}

\newtheorem{definition}{Definition}
\newtheorem{proposition}{Proposition}
\newtheorem{remark}{Remark}

\newcommand{\PP}{\mathbb{P}}
\newcommand{\EE}{\mathbb{E}}
\newcommand{\defeq}{\stackrel{\mathrm{def}}{=}}
\newcommand{\MaxL}{\mathcal{L}}

\title{Spectral Delegation for Anonymous DHT Lookups\\
\large From mixing certificates to auditable per-step budgets}
\author{}
\date{March 16, 2026 (archive v394)}

\begin{document}
\maketitle

\begin{abstract}
Random-walk delegation is a standard building block for ``anonymous'' DHT lookups:
a querier hands work to a relay chosen via a short walk on an overlay, and the relay performs the visible step.
The published spectral anonymity note gives sharp distinguishing bounds for delegation in (possibly directed) overlays.
This paper extracts only the \emph{spectral-witness $\Rightarrow$ auditable per-step budget} corollary: a compact witness (e.g.\ a second-singular-value bound)
can be turned into a checked per-step MaxL (bits) cap; via the endpoint bridge this immediately yields an odds-inflation bound and identification sizing rules.\cite{series:endpointbridge}
The paper stops there.  Notarized evaluation bundles, ABOM/OINL claim surfaces, and release-lineage / signed-receipt packaging belong to companion papers.\cite{series:notcert,series:abomoinl,series:releaseA}
Mechanism papers should cite this note (and the published spectral note) rather than re-deriving spectral-to-budget translations.
\end{abstract}

\paragraph{Scope and non-redundancy.}
We only do the ``spectral $\Rightarrow$ per-step budget'' translation.
Receipt schema, replay hooks, adaptive horizon accounting, notarized evaluator bundles, and release packaging are handled elsewhere.\cite{series:receiptitems,series:condcert,series:accountingrulebook,series:notcert,series:abomoinl,series:releaseA}
Systems and related-work context for anonymous DHT delegation is indexed in Synthesis~7.\cite{series:relatedworkmap}

\paragraph{Canonical backbone.}
The full spectral anonymity result and optimal distinguishing bounds are owned by the published mathematics note:
\path{[[2026.01.23 - Mathematics: Spectral Anonymity and Optimal Distinguishing Bounds for Random-Walk Delegation in (Possibly Directed) Overlays]]}.

\paragraph{Units.}
Budgets are measured in bits: we use $\log_2$ and $2^{\epsilon}$ for odds-inflation / pointwise-maximal-leakage style caps.\cite{series:notation}

\paragraph{Minimal exported object.}
The stable output of this paper is the tiny witness-to-budget object
\begin{quote}\small
\begin{tabular}{@{}l@{}}
\path|witness_mode, kernel_digest, prior_scope,|\\
\path|  n, t, sigma2_upper or chi2_bound,|\\
\path|  pi_min (optional), delta, derived_budget_bits.|
\end{tabular}
\end{quote}
where \texttt{witness\_mode} is either \texttt{spectral-uniform} (for \cref{prop:spectral-to-maxl}) or \texttt{chi2-general} (for \cref{prop:chi2-to-maxl}).
This paper does \emph{not} own the enclosing notarized bundle, ABOM summary, or cross-release lineage declaration; it only fixes how the numerical cap is recomputed from the compact witness.\cite{series:notcert,series:abomoinl,series:releaseA}

\paragraph{Later-terse-paper citation posture.}
Later terse papers should cite this note only when the live issue is the witness-to-budget translation itself: which kernel digest, which \texttt{witness\_mode}, and which audited per-step cap were used. If the live issue is the theorem root, stop at the published Mathematics spectral note; if the live issue is the Regime~II certificate shape, stop at Synthesis~21; if the live issue is replaying a delivered receipt bundle, stop at Synthesis~20; and if the live issue is evaluator notarization or release lineage, stop at Eval~1 or Release~A respectively rather than widening this note into a generic delegation dossier.\cite{series:spectral,series:condcert,series:auditorreplay,series:notcert,series:releaseA}

\section{Delegation model and spectral witness}
\label{sec:model}

We adopt the delegation setting of the published spectral anonymity paper.\cite{series:spectral}
A start node $S$ (the querier, or the querier's entry relay) initiates a length-$t$ delegation walk on an $n$-node overlay.
The observed delegate $D$ is the terminal node of the walk (or, more generally, the node that performs a visible lookup step).

The published note defines a discriminant operator for possibly directed kernels and shows that distinguishing power contracts at rate controlled by a
second singular value $\sigma_2\in[0,1)$ of that operator (or a conservative certificate upper-bounding it).\cite{series:spectral}
For this paper, the \emph{witness} is simply: the overlay size $n$, the walk length $t$, the kernel identifier (a digest), and an auditor-checkable bound $\sigma_2$.

\begin{remark}[Why a spectral witness fits receipts]
A receipt cannot enumerate a huge conditioning space (all possible route histories), but it \emph{can} commit to a small, digest-bound kernel and a small numerical witness.
Regime~II conditional certificates (Synthesis~21) are designed to carry exactly that shape: a conditioning interface identifier plus a compact evidence object.\cite{series:condcert}
\end{remark}


\section{A conservative spectral-to-budget translation}
\label{sec:translation}

The published spectral note provides (among stronger results) a conservative bound on expected $\chi^2$ leakage of the posterior on $S$ given $D$.\cite{series:spectral}
We record a short corollary that yields an auditable MaxL cap in bits with a tail parameter $\delta$.

\begin{proposition}[Spectral contraction $\Rightarrow$ a high-probability per-step MaxL bound (uniform prior)]
\label{prop:spectral-to-maxl}
Assume the stationary prior $\pi$ is uniform on $n$ nodes.
Let $\sigma_2$ be an upper bound on the second singular value of the discriminant operator for the delegation kernel (as defined in the published spectral note).\cite{series:spectral}
Let $\mu_D(\cdot)\defeq \PP[S=\cdot\mid D]$ be the posterior on $S$ given the observed delegate $D$.
Then for any $\delta\in(0,1)$, with probability at least $1-\delta$ over $D\sim\pi$,
\[
\MaxL(S\!\to\!D)
\;\le\;
\log_2\!\Big(1+\sqrt{n(n-1)}\,\sigma_2^{t}/\sqrt{\delta}\Big)\ \text{bits}.
\]
\end{proposition}

\begin{proof}[Proof sketch]
The published spectral note gives
$\EE_{D\sim\pi}\!\left[\chi^2(\mu_D\|\pi)\right]\le (n-1)\sigma_2^{2t}$.\cite{series:spectral}
Markov's inequality yields $\chi^2(\mu_D\|\pi)\le (n-1)\sigma_2^{2t}/\delta$ with probability at least $1-\delta$.
For uniform $\pi$, $\chi^2(\mu\|\pi)=n\sum_s(\mu(s)-1/n)^2$, so
$\|\mu-\pi\|_2 = \sqrt{\chi^2(\mu\|\pi)/n}$ and $\|\mu-\pi\|_\infty\le \|\mu-\pi\|_2$.
Thus $\max_s \mu(s)\le 1/n+\sqrt{\chi^2(\mu\|\pi)/n}$, implying
$\max_s \mu(s)/\pi(s)\le 1+\sqrt{n\,\chi^2(\mu\|\pi)}$.
Taking $\log_2$ gives the stated bound.
\end{proof}

\begin{remark}[What this bound is (and is not)]
This corollary is intentionally conservative and assumes a uniform stationary prior.
It is nevertheless useful as a receipt-side \emph{auditable} cap: the parameters $(n,t,\sigma_2,\delta)$ are small, and the auditor can recompute the right-hand side.
Sharper and more general statements (non-uniform priors, directed kernels, optimal tests) live in the published spectral note.\cite{series:spectral}
\end{remark}

\begin{proposition}[From a $\chi^2$ witness to a MaxL cap (non-uniform priors)]
\label{prop:chi2-to-maxl}
Let $\pi$ be any prior on $S$ with $\pi_{\min}\defeq \min_s \pi(s)>0$.
For any posterior $\mu$ on $S$,
\[
\max_s \frac{\mu(s)}{\pi(s)}\;\le\;1+\sqrt{\chi^2(\mu\|\pi)/\pi_{\min}}.
\]
Consequently, if an evidence object implies that with probability at least $1-\delta$ (over $D$) we have $\chi^2(\mu_D\|\pi)\le B$, then with the same probability
\[
\MaxL(S\!\to\!D)\;\le\;\log_2\!\Big(1+\sqrt{B/\pi_{\min}}\Big)\ \text{bits}.
\]
\end{proposition}

\begin{proof}[Proof sketch]
By definition, $\chi^2(\mu\|\pi)=\sum_s (\mu(s)-\pi(s))^2/\pi(s)$, so for each $s$,
$(\mu(s)-\pi(s))^2/\pi(s)\le \chi^2(\mu\|\pi)$ and hence
$|\mu(s)-\pi(s)|\le \sqrt{\pi(s)\,\chi^2(\mu\|\pi)}$.
Therefore
$\mu(s)/\pi(s)\le 1+\sqrt{\chi^2(\mu\|\pi)/\pi(s)}\le 1+\sqrt{\chi^2(\mu\|\pi)/\pi_{\min}}$.
Take $\log_2$.
\end{proof}


\section{Handoff to certificates and release objects}
\label{sec:receipt-use}

A deployment that uses delegation to hide a querier's identity (or entry relay) can publish a Regime~II per-step conditional certificate whose evidence object is:

\begin{quote}\small
\begin{tabular}{@{}l@{}}
\path|kernel_digest, n, t, sigma2_upper, delta|\\
\path|or|\\
\path|kernel_digest, prior_scope, chi2_bound, delta.|
\end{tabular}
\end{quote}

The replay plan commits to the kernel digest and the numerical witness.
An auditor recomputes the per-step cap $b_t$ from \cref{prop:spectral-to-maxl} or \cref{prop:chi2-to-maxl} and records only that checked conditional budget for the step.
Adaptive horizon accounting then composes the checked per-step caps additively (Synthesis~19), without assuming independence.\cite{series:accountingrulebook,series:condcert}

\paragraph{Earliest sufficient stop.}
For a conservative audit spine, the delegation argument can stop after three steps:
(1) cite the published spectral mathematics note for the kernel family;\cite{series:spectral}
(2) attach the compact witness from this paper and recompute the per-step budget cap; and
(3) hand that checked cap either to the notarized evaluation bundle (Eval~1) or to the ABOM/OINL + Release~A publication surface.\cite{series:notcert,series:abomoinl,series:releaseA}
Nothing in this paper requires re-explaining those later layers.

\paragraph{Serialization rule.}
To keep certificate-facing companions from re-specifying this translation, the replay object should use exactly one of the following two minimal modes:
\begin{enumerate}[leftmargin=1.4em]
\item \textbf{\texttt{spectral-uniform}:} \texttt{kernel\_digest, n, t, sigma2\_upper, delta} with the auditor recomputing
\[
 b_t = \log_2\!\Big(1+\sqrt{n(n-1)}\,\sigma_2^t/\sqrt{\delta}\Big)
\]
from \cref{prop:spectral-to-maxl};
\item \textbf{\texttt{chi2-general}:} \texttt{\{ kernel\_digest, prior\_scope, chi2\_bound, pi\_min, delta \}} with the auditor recomputing
\[
 b_t = \log_2\!\Big(1+\sqrt{\chi^2\_\text{bound}/\pi_{\min}}\Big)
\]
from \cref{prop:chi2-to-maxl}.
\end{enumerate}
Any richer theorem-side object (directed-kernel witness, time-varying schedule digest, class-prior certificate) should be reduced to one of these two replay modes before it reaches the receipt layer.
That keeps this paper small while still making the exported witness object explicit enough for Regime~II certificates and Eval~1 bundles to reuse without reinterpretation.\cite{series:condcert,series:notcert}

\paragraph{A minimal public spectral-witness card.}
\begin{center}\small
\begin{tabular}{@{}ll@{}}
Claim scope & Delegation step from pinned kernel digest \texttt{ker:7c91\ldots} \\
Replay mode & \texttt{spectral-uniform} \\
Witness & $n=1024$, $t=8$, $\sigma_2^{\mathrm{upper}}=0.25$, $\delta=10^{-2}$ \\
Derived per-step budget & $b_t=\log_2(1+\sqrt{1024\cdot 1023}\cdot 0.25^8/\sqrt{10^{-2}})\approx 0.209$ bits \\
Interpretation & per-step odds multiplier at most $2^{b_t}\approx 1.156$ with probability at least $0.99$ \\
Downstream handoff & checked cap only; notarization / ABOM / release lineage omitted here \\
\end{tabular}
\end{center}
This is the smallest public packet later terse papers need when the live question is simply which compact spectral witness was checked and what per-step budget it yielded. Anything richer belongs either in the theorem root or in the downstream certificate / release companions already named above.

\paragraph{Arithmetic drill.}
The arithmetic is intentionally mechanical: $\sqrt{1024\cdot 1023}\approx 1023.50$, $0.25^8=1/65536$, and so the inner factor is about $0.1562$; hence $b_t=\log_2(1+0.1562)\approx 0.209$ bits. An auditor does not need the walk transcript or a deployment diary here---only the pinned kernel digest, the declared witness, and the deterministic recomputation rule from \cref{prop:spectral-to-maxl}.

\paragraph{Where to put the witness.}
This paper does not add new public-facing fields.
The evidence object lives behind the existing replay hook (plan identifier + plan-spec digest), while Eval~1 owns the notarized evaluator bundle and Release~A owns the release-facing lineage / signed-receipt layer.\cite{series:auditorreplay,series:notcert,series:releaseA}

\paragraph{Non-uniform priors (class budgets).}
If the operator budgets ``user within a class'' (non-uniform $\pi$), keep the outward claim surface unchanged: include a prior certificate parameter such as $\pi_{\min}$ and use \cref{prop:chi2-to-maxl} to translate any $\chi^2$ witness into a MaxL cap.
The published spectral note contains the sharper (and directed-kernel) statements that supply such witnesses; this paper only packages the final auditor-side conversion.\cite{series:spectral}

\section{Research warranted}
The most useful refinements are interface refinements, not longer papers:
(i) non-uniform priors and class-based priors (so the budget targets ``user within a set'');
(ii) directed overlays and time-varying kernels (so the kernel digest is a schedule rather than a single matrix); and
(iii) a tiny standardized field vocabulary for spectral witnesses so the same witness can flow unchanged through Eval~1 and Release~A.
Both (i) and (ii) are treated in the published spectral note at the level of optimal distinguishing bounds;\cite{series:spectral}
the remaining work is to package those bounds into compact, auditor-friendly evidence objects rather than to enlarge the theorem surface.

\begin{thebibliography}{99}\small

\bibitem{series:notation}
Anonymous.
\newblock Notation and Model Conventions (bits, logs, identifiers).
\newblock Synthesis~8, The-One-True-Archive (2026).


\bibitem{series:endpointbridge}
Anonymous.
\newblock Endpoint Metrics Bridge: Odds Inflation, PML/MaxL, and Identification Sizing.
\newblock Synthesis~5, The-One-True-Archive (2026).

\bibitem{series:relatedworkmap}
Anonymous.
\newblock Related Work Map for Anonymous Overlays and DHT Privacy.
\newblock Synthesis~7, The-One-True-Archive (2026).

\bibitem{series:spectral}
Anonymous.
\newblock Spectral Anonymity and Optimal Distinguishing Bounds for Random-Walk Delegation in (Possibly Directed) Overlays.
\newblock Published note, [[2026.01.23 - Mathematics: Spectral Anonymity and Optimal Distinguishing Bounds for Random-Walk Delegation in (Possibly Directed) Overlays]], 2026. (This archive.)

\bibitem{series:notcert}
Anonymous.
\newblock Notarized Anonymity Certificates Under Retries: Attempt Logs, Spending Rules, and Verifier Bundles.
\newblock Evaluation~1, The-One-True-Archive (2026).

\bibitem{series:abomoinl}
Anonymous.
\newblock ABOM and OINL for Anonymity Claims: Digest-Bound Objects and Reader-Facing Labels.
\newblock Anonymity~C, The-One-True-Archive (2026).

\bibitem{series:releaseA}
Anonymous.
\newblock Privacy as a Release Property: Lineage, Receipts, and Anti-Equivocation for Anonymous DHT Claims.
\newblock Release~A, The-One-True-Archive (2026).

\bibitem{series:receiptitems}
Anonymous.
\newblock Receipt Line Items for Anonymity Claims: A Minimal Tuple Schema for Published Budgets.
\newblock Synthesis~12, The-One-True-Archive (2026).

\bibitem{series:condcert}
Anonymous.
\newblock Conditional Per-Step Budget Certificates: A Receipt-Grade Interface for Adaptive Horizon Claims.
\newblock Synthesis~21, The-One-True-Archive (2026).

\bibitem{series:accountingrulebook}
Anonymous.
\newblock Receipt Accounting Rulebook: Repeats, Retries, and State.
\newblock Synthesis~19, The-One-True-Archive (2026).

\bibitem{series:auditorreplay}
Anonymous.
\newblock Auditor Replay Recipe: Mechanical Checks for Receipts, Plans, and Evidence Objects.
\newblock Synthesis~20, The-One-True-Archive (2026).

\end{thebibliography}

\end{document}
